\scanpage{37}
Dans la situation générale des ext. pour $G$ correspondant
à une extension

\diagram{\begin{tikzcd}[ampersand replacement=\&,column sep=small]
1\arrow[r,dash]\&\pi_1\arrow[r,dash]\&G_1\arrow[r,dash,"d"]\&G_0\arrow[r,dash]\arrow[d,dashed,"\theta"]\&\pi_0\arrow[r,dash]\&1\\
\&\&\&\operatorname{Aut}(G_1,\pi_1)\&\&
\end{tikzcd}}

on a (pour le $G$ donné) un hom
\[
\operatorname{Cent}(G_0)\cap\operatorname{Ker}\theta\xrightarrow{\varphi_G}Z^1(\pi_0,\pi_1)
\]
dont le noyau est l'image de $Z(G)$ dans $G_0$.

\margin{$\overbrace{\underbrace{\pi_0,\operatorname{Im}d}_{G_0},\pi_1}^{G}$}

Dans le cas qui nous intéresse
\[
\begin{aligned}
\relax[\ G_1&=D'=L_1.D_1=\operatorname{Cent}G.\operatorname{Der}G,\quad
G_0=L'=G/L_1\cap D_1\\
&=G/\operatorname{Cent}G\cap\operatorname{Der}G\quad ---\ ]
\end{aligned}
\]
cela \deleted{\unsure{aussi}} s'exprime par
\[
\operatorname{Cent}(D_0)\cap\operatorname{Ker}\theta\xrightarrow{\varphi_G}\operatorname{Hom}(\pi_0,\pi_1)
\]
\inserted{(c'est ce qui exprime que $L_0\longleftrightarrow\operatorname{Cent}G$ est iso)}
dont le noyau doit être nul; ceci montre
donc qu'on doit avoir (compte tenu
que $\pi_1$ est quasi-unipot.)
\[
\boxed{\operatorname{Cent}D_0\cap\operatorname{Ker}\theta\text{ est quasi-unipotent}}
\]
\inserted{$\simeq[\operatorname{Cent}(\operatorname{D\acute{e}r}(G))/\operatorname{Cent}(G)]\cap\operatorname{Cent}D'_0$}
\margin{$D_0\simeq G/\unsure{(\operatorname{Cent}G)}$}

En fait, on doit avoir mieux, comme
\unsure{à savoir} \uline{$\operatorname{Cent}(D_0)$ unipotent}

\diagram{\begin{tikzcd}[ampersand replacement=\&,column sep=small,row sep=large]
\&\&L_1\arrow[dl,dash,hook',"A"']\arrow[dr,dash,hook',"H"]\arrow[drr,bend left=40,dash,"D_0"]\&\&\\
1\arrow[r,dash,hook]\&\pi_1=L_1\cap D_1\arrow[rr,dash,hook,"A\times H"]\arrow[dr,dash,hook,"H"']\&\&D'\arrow[r,dash,hook]\&G\\
\&\&D_1\arrow[ur,dash,hook,"A"']\arrow[urr,bend right=40,dash,"L_0"']\&\&
\end{tikzcd}}
\margin{\unsure{donc}}
\mspage{36}
